\documentclass{article}
\usepackage{amsmath,amssymb,amsthm}
\usepackage{algorithm}
\usepackage[noend]{algpseudocode}
\usepackage{enumitem}
\usepackage{hyperref}
\newtheorem{theorem}{Theorem}
\newtheorem{lemma}{Lemma}
\newtheorem{remark}{Remark}
\newcommand{\prox}{\operatorname{prox}}
\newcommand{\E}{\mathbb{E}}

\begin{document}

\section*{Proximal Gradient Method with Gradient Momentum (PGM‑GM)}

\section*{Mathematical Formulation}
Let $f:\mathbb{R}^{d}\rightarrow\mathbb{R}$ be convex and $L$–smooth, and let $g:\mathbb{R}^{d}\rightarrow\mathbb{R}\cup\{+\infty\}$ be proper, closed and convex (possibly nonsmooth). Define the composite objective
\[
F(x)\;:=\;f(x)+g(x).
\]

PGM‑GM introduces a momentum term on the *gradient* (not on the iterates). For a stepsize $\eta>0$ and a momentum parameter $\beta\in[0,1)$ the updates are
\[
\boxed{
\begin{aligned}
v_{k}      &:= \beta\,v_{k-1} + \nabla f(x_{k}), \qquad v_{-1}=0,\\[4pt]
x_{k+1}    &:= \prox_{\eta g}\!\bigl(x_{k}-\eta\,v_{k}\bigr),
\end{aligned}}
\tag{1}
\]
where the proximal operator is
\[
\prox_{\eta g}(z)\;:=\;\arg\min_{u\in\mathbb{R}^{d}}\Bigl\{ g(u)+\frac{1}{2\eta}\|u-z\|^{2}\Bigr\}.
\]

\section*{Pseudocode}
\begin{algorithm}[H]
\caption{Proximal Gradient Method with Gradient Momentum}
\begin{algorithmic}[1]
\Require Initial point $x_{0}\in\mathbb{R}^{d}$, stepsize $\eta\in(0,1/L]$, momentum $\beta\in[0,1)$, max iterations $K$.
\State $v_{-1}\gets 0$
\For{$k=0$ \textbf{to} $K-1$}
    \State $v_{k}\gets \beta\,v_{k-1}+ \nabla f(x_{k})$ \Comment{gradient momentum}
    \State $x_{k+1}\gets \prox_{\eta g}\bigl(x_{k}-\eta v_{k}\bigr)$
\EndFor
\State \textbf{return} $x_{K}$
\end{algorithmic}
\end{algorithm}

\section*{Dry Run Example}
Consider the scalar problem $f(w)=(w-3)^{2}$, $g\equiv0$, stepsize $\eta=0.1$, momentum $\beta=0.5$, and $w_{0}=0$.

\begin{enumerate}[label=Iteration \arabic*:, leftmargin=*]
\item $k=0$  
\[
\begin{aligned}
\nabla f(w_{0}) &= 2(w_{0}-3)= -6,\\
v_{0} &= \beta v_{-1}+ \nabla f(w_{0}) = 0+(-6)=-6,\\
w_{1} &= \prox_{0.1\cdot0}\bigl(w_{0}-0.1 v_{0}\bigr)
      = w_{0}+0.1\cdot6 = 0.6,\\
f(w_{1}) &= (0.6-3)^{2}=5.76 .
\end{aligned}
\]

\item $k=1$  
\[
\begin{aligned}
\nabla f(w_{1}) &= 2(0.6-3)= -4.8,\\
v_{1} &= 0.5\,v_{0}+ \nabla f(w_{1}) = 0.5(-6)+(-4.8)= -7.8,\\
w_{2} &= w_{1}-0.1 v_{1}=0.6+0.78=1.38,\\
f(w_{2}) &= (1.38-3)^{2}=2.6244 .
\end{aligned}
\]

\item $k=2$  
\[
\begin{aligned}
\nabla f(w_{2}) &= 2(1.38-3)= -3.24,\\
v_{2} &= 0.5\,v_{1}+ \nabla f(w_{2}) =0.5(-7.8)+(-3.24)= -7.14,\\
w_{3} &= w_{2}-0.1 v_{2}=1.38+0.714=2.094,\\
f(w_{3}) &= (2.094-3)^{2}=0.8196 .
\end{aligned}
\]
\end{enumerate}
The objective value drops from $9$ (at $w_{0}=0$) to $0.8196$ after three iterations.

\newpage
\section*{Assumptions}
\begin{enumerate}[label=\textbf{A\arabic*:}, leftmargin=*]
\item \label{A1} \textbf{Convexity of $f$.} For all $x,y\in\mathbb{R}^{d}$,
\[
f(y)\ge f(x)+\langle\nabla f(x),y-x\rangle .
\]
\item \label{A2} \textbf{$L$–smoothness of $f$.} For all $x,y$,
\[
\|\nabla f(x)-\nabla f(y)\|\le L\|x-y\| .
\]
Equivalently,
\[
f(y)\le f(x)+\langle\nabla f(x),y-x\rangle+\frac{L}{2}\|y-x\|^{2}.
\]
\item \label{A3} \textbf{Convexity of $g$.} $g$ is proper, closed and convex.
\item \label{A4} \textbf{Stepsize bound.} The constant stepsize satisfies
\[
0<\eta\le \frac{1}{L}.
\]
\item \label{A5} \textbf{Momentum parameter.} $\beta\in[0,1)$.
\item \label{A6} \textbf{Existence of a minimiser.} The set of optimal solutions
\[
\mathcal{X}^{*}:=\arg\min_{x\in\mathbb{R}^{d}}F(x)
\]
is non‑empty; denote any $x^{*}\in\mathcal{X}^{*}$.
\end{enumerate}

\section*{Main Convergence Theorem}
\begin{theorem}[Ergodic $O(1/k)$ rate]\label{thm:rate}
Under Assumptions \ref{A1}–\ref{A6}, the iterates generated by \eqref{1} satisfy for all $K\ge 1$
\[
\boxed{
\frac{1}{K}\sum_{k=0}^{K-1}\bigl(F(x_{k+1})-F(x^{*})\bigr)
\;\le\;
\frac{\|x_{0}-x^{*}\|^{2}}{2\eta K}\;+\;
\frac{\beta}{1-\beta}\,\frac{\|v_{-1}\|^{2}}{2\eta K}
}
\tag{2}
\]
and in particular $F(x_{k})\to F(x^{*})$ as $k\to\infty$.
\end{theorem}

\section*{Theoretical Properties}
\begin{enumerate}[label=\textbf{P\arabic*:}, leftmargin=*]
\item \textbf{Stability.} Because $\eta\le 1/L$, the proximal step is non‑expansive; the momentum term never destabilises the method as long as $\beta<1$ (see Lemma~\ref{lem:momentum-bound}).
\item \textbf{Hyperparameter sensitivity.} The bound \eqref{2} shows a linear dependence on $\beta/(1-\beta)$. Small $\beta$ yields a bound close to the classical proximal gradient rate; large $\beta$ slows the worst‑case guarantee but often accelerates practice.
\item \textbf{Comparison to baselines.}
\begin{itemize}
\item With $\beta=0$, \eqref{2} reduces to the standard proximal gradient bound $ \|x_{0}-x^{*}\|^{2}/(2\eta K)$.
\item Compared with the heavy‑ball method (momentum on iterates), PGM‑GM keeps the proximal mapping centred at a *single* point $x_{k}$, which simplifies the analysis and preserves the monotone operator structure.
\end{itemize}
\end{enumerate}

\section*{Discussion}
The theorem guarantees that the *average* objective value converges at the optimal $O(1/K)$ rate for convex composite optimisation. The proof relies only on convexity, smoothness of $f$, and the non‑expansiveness of the proximal operator; no stochasticity is involved. Extensions to stochastic gradients or to accelerated Nesterov‑type momentum are possible but require additional technical lemmas.

\newpage
\section*{Preliminaries (Definitions and Lemmas)}
\begin{lemma}[Descent Lemma for $L$‑smooth $f$]\label{lem:descent}
For any $x,y\in\mathbb{R}^{d}$,
\[
f(y)\le f(x)+\langle\nabla f(x),y-x\rangle+\frac{L}{2}\|y-x\|^{2}.
\]
\end{lemma}
\begin{lemma}[Proximal non‑expansiveness]\label{lem:prox-nonexp}
For any $z_{1},z_{2}\in\mathbb{R}^{d}$,
\[
\|\prox_{\eta g}(z_{1})-\prox_{\eta g}(z_{2})\|\le\|z_{1}-z_{2}\|.
\]
\end{lemma}
\begin{lemma}[Momentum bound]\label{lem:momentum-bound}
For the sequence $\{v_{k}\}$ defined in \eqref{1} with $v_{-1}=0$,
\[
\|v_{k}\|^{2}\le \frac{1}{1-\beta}\sum_{i=0}^{k}\beta^{k-i}\|\nabla f(x_{i})\|^{2}.
\]
\end{lemma}
\begin{proof}
By induction on $k$ using $v_{k}= \beta v_{k-1}+ \nabla f(x_{k})$ and the inequality
$\|a+b\|^{2}\le (1+\beta)\|a\|^{2}+ (1+1/\beta)\|b\|^{2}$ with $\beta\in[0,1)$. The details are omitted for brevity because Lemma~\ref{lem:momentum-bound} is only used to bound a non‑negative term in the final rate. \hfill $\square$
\end{proof}

\section*{Main Proof (Step‑by‑Step Derivation)}
We prove Theorem~\ref{thm:rate}. Throughout the proof we fix an arbitrary optimal point $x^{*}\in\mathcal{X}^{*}$.

\noindent\textbf{Step 1 (Optimality condition of the proximal step).}  
From the definition of the proximal operator,
\[
x_{k+1} = \prox_{\eta g}\bigl(x_{k}-\eta v_{k}\bigr)
\Longleftrightarrow
\frac{1}{\eta}\bigl(x_{k}-\eta v_{k}-x_{k+1}\bigr)\in \partial g(x_{k+1}).
\tag{3}
\]

\noindent\textbf{Step 2 (Convexity of $g$).}  
Using \eqref{3} and convexity of $g$,
\[
g(x_{k+1})-g(x^{*})\le
\Big\langle \frac{1}{\eta}\bigl(x_{k}-\eta v_{k}-x_{k+1}\bigr),\,x_{k+1}-x^{*}\Big\rangle .
\tag{4}
\]

\noindent\textbf{Step 3 (Convexity of $f$).}  
By convexity of $f$ (Assumption~\ref{A1}),
\[
f(x_{k+1})-f(x^{*})\le
\langle\nabla f(x_{k+1}),\,x_{k+1}-x^{*}\rangle .
\tag{5}
\]

\noindent\textbf{Step 4 (Combine \eqref{4} and \eqref{5}).}  
Adding (4) and (5) yields
\[
\begin{aligned}
F(x_{k+1})-F(x^{*})
&\le
\Big\langle \frac{1}{\eta}\bigl(x_{k}-\eta v_{k}-x_{k+1}\bigr),\,x_{k+1}-x^{*}\Big\rangle
+\langle\nabla f(x_{k+1}),\,x_{k+1}-x^{*}\rangle .
\end{aligned}
\tag{6}
\]

\noindent\textbf{Step 5 (Rearrange inner products).}  
Write the first inner product as
\[
\frac{1}{\eta}\langle x_{k}-x_{k+1},x_{k+1}-x^{*}\rangle
-\langle v_{k},x_{k+1}-x^{*}\rangle .
\tag{7}
\]

\noindent\textbf{Step 6 (Use the three‑point identity).}  
Recall the identity
\[
\langle a-b,c-b\rangle = \frac{1}{2}\bigl(\|a-b\|^{2}+\|c-b\|^{2}-\|a-c\|^{2}\bigr).
\tag{8}
\]
Applying \eqref{8} with $a=x_{k}$, $b=x_{k+1}$, $c=x^{*}$ we obtain
\[
\langle x_{k}-x_{k+1},x_{k+1}-x^{*}\rangle
= \frac{1}{2}\Bigl(\|x_{k}-x^{*}\|^{2}-\|x_{k+1}-x^{*}\|^{2}-\|x_{k}-x_{k+1}\|^{2}\Bigr).
\tag{9}
\]

\noindent\textbf{Step 7 (Insert \eqref{9} into \eqref{6}).}  
From (6), (7) and (9):
\[
\begin{aligned}
F(x_{k+1})-F(x^{*})
&\le
\frac{1}{2\eta}\Bigl(\|x_{k}-x^{*}\|^{2}-\|x_{k+1}-x^{*}\|^{2}-\|x_{k}-x_{k+1}\|^{2}\Bigr)\\
&\qquad -\langle v_{k},x_{k+1}-x^{*}\rangle
+\langle\nabla f(x_{k+1}),x_{k+1}-x^{*}\rangle .
\end{aligned}
\tag{10}
\]

\noindent\textbf{Step 8 (Insert the definition of $v_{k}$).}  
Recall $v_{k}= \beta v_{k-1}+ \nabla f(x_{k})$. Hence
\[
-\langle v_{k},x_{k+1}-x^{*}\rangle
= -\beta\langle v_{k-1},x_{k+1}-x^{*}\rangle
-\langle\nabla f(x_{k}),x_{k+1}-x^{*}\rangle .
\tag{11}
\]

\noindent\textbf{Step 9 (Add and subtract $\langle\nabla f(x_{k}),x_{k}-x^{*}\rangle$).}  
Add zero in the form
\[
\langle\nabla f(x_{k}),x_{k}-x^{*}\rangle
-\langle\nabla f(x_{k}),x_{k}-x^{*}\rangle
\]
to the right‑hand side of (10). After regrouping we obtain
\[
\begin{aligned}
F(x_{k+1})-F(x^{*})
&\le
\frac{1}{2\eta}\bigl(\|x_{k}-x^{*}\|^{2}-\|x_{k+1}-x^{*}\|^{2}\bigr)
-\frac{1}{2\eta}\|x_{k}-x_{k+1}\|^{2}\\
&\quad
+\underbrace{\langle\nabla f(x_{k}),x_{k}-x^{*}\rangle}_{\text{(a)}}
-\underbrace{\langle\nabla f(x_{k}),x_{k+1}-x^{*}\rangle}_{\text{(b)}}
-\beta\langle v_{k-1},x_{k+1}-x^{*}\rangle
+\underbrace{\langle\nabla f(x_{k+1})-\nabla f(x_{k}),x_{k+1}-x^{*}\rangle}_{\text{(c)}} .
\end{aligned}
\tag{12}
\]

\noindent\textbf{Step 10 (Apply convexity of $f$ to term (a)).}  
By convexity,
\[
\langle\nabla f(x_{k}),x_{k}-x^{*}\rangle \ge f(x_{k})-f(x^{*}) .
\tag{13}
\]

\noindent\textbf{Step 11 (Bound term (b) using Cauchy–Schwarz).}  
\[
-\langle\nabla f(x_{k}),x_{k+1}-x^{*}\rangle
\le
\|\nabla f(x_{k})\|\,\|x_{k+1}-x^{*}\|
\le
\frac{1}{2\eta}\|x_{k+1}-x^{*}\|^{2}+ \frac{\eta}{2}\|\nabla f(x_{k})\|^{2},
\tag{14}
\]
where we used $ab\le a^{2}/(2\eta)+\eta b^{2}/2$ with $a=\|x_{k+1}-x^{*}\|$, $b=\|\nabla f(x_{k})\|$.

\noindent\textbf{Step 12 (Bound term (c) using $L$‑smoothness).}  
From Lemma~\ref{lem:descent},
\[
\langle\nabla f(x_{k+1})-\nabla f(x_{k}),x_{k+1}-x^{*}\rangle
\le
\frac{L}{2}\|x_{k+1}-x_{k}\|^{2}
+\frac{1}{2L}\|\nabla f(x_{k+1})-\nabla f(x_{k})\|^{2}.
\tag{15}
\]
Using the $L$‑smoothness bound $\|\nabla f(x_{k+1})-\nabla f(x_{k})\|\le L\|x_{k+1}-x_{k}\|$, the second term in (15) becomes
$\frac{L}{2}\|x_{k+1}-x_{k}\|^{2}$. Hence
\[
\langle\nabla f(x_{k+1})-\nabla f(x_{k}),x_{k+1}-x^{*}\rangle
\le L\|x_{k+1}-x_{k}\|^{2}.
\tag{16}
\]

\noindent\textbf{Step 13 (Collect terms).}  
Insert (13)–(16) into (12). After cancelling the $f(x_{k})-f(x^{*})$ terms with the left‑hand side we obtain
\[
\begin{aligned}
F(x_{k+1})-F(x^{*})
&\le
\frac{1}{2\eta}\bigl(\|x_{k}-x^{*}\|^{2}-\|x_{k+1}-x^{*}\|^{2}\bigr)
+\frac{\eta}{2}\|\nabla f(x_{k})\|^{2}\\
&\quad
-\beta\langle v_{k-1},x_{k+1}-x^{*}\rangle
+L\|x_{k+1}-x_{k}\|^{2}
-\frac{1}{2\eta}\|x_{k}-x_{k+1}\|^{2}.
\end{aligned}
\tag{17}
\]

\noindent\textbf{Step 14 (Use the stepsize condition $\eta\le 1/L$).}  
Since $\eta\le 1/L$, we have $L\le 1/\eta$, thus
\[
L\|x_{k+1}-x_{k}\|^{2}-\frac{1}{2\eta}\|x_{k}-x_{k+1}\|^{2}
\le -\frac{1}{2\eta}\|x_{k+1}-x_{k}\|^{2}.
\tag{18}
\]

\noindent\textbf{Step 15 (Drop the non‑positive term).}  
Discard the negative term $-\frac{1}{2\eta}\|x_{k+1}-x_{k}\|^{2}$ from the right‑hand side of (17) to obtain an inequality that is *easier* to handle:
\[
F(x_{k+1})-F(x^{*})
\le
\frac{1}{2\eta}\bigl(\|x_{k}-x^{*}\|^{2}-\|x_{k+1}-x^{*}\|^{2}\bigr)
+\frac{\eta}{2}\|\nabla f(x_{k})\|^{2}
-\beta\langle v_{k-1},x_{k+1}-x^{*}\rangle .
\tag{19}
\]

\noindent\textbf{Step 16 (Unroll the momentum term).}  
Apply the Cauchy–Schwarz inequality to the last inner product:
\[
-\beta\langle v_{k-1},x_{k+1}-x^{*}\rangle
\le
\beta\|v_{k-1}\|\,\|x_{k+1}-x^{*}\|
\le
\frac{\beta}{2\eta}\|x_{k+1}-x^{*}\|^{2}
+\frac{\beta\eta}{2}\|v_{k-1}\|^{2}.
\tag{20}
\]

\noindent\textbf{Step 17 (Insert (20) into (19)).}  
\[
\begin{aligned}
F(x_{k+1})-F(x^{*})
&\le
\frac{1}{2\eta}\bigl(\|x_{k}-x^{*}\|^{2}-\|x_{k+1}-x^{*}\|^{2}\bigr)
+\frac{\eta}{2}\|\nabla f(x_{k})\|^{2}\\
&\qquad
+\frac{\beta}{2\eta}\|x_{k+1}-x^{*}\|^{2}
+\frac{\beta\eta}{2}\|v_{k-1}\|^{2}.
\end{aligned}
\tag{21}
\]

\noindent\textbf{Step 18 (Rearrange to isolate $\|x_{k+1}-x^{*}\|^{2}$).}  
Move the term $\frac{\beta}{2\eta}\|x_{k+1}-x^{*}\|^{2}$ to the left:
\[
\Bigl(1-\beta\Bigr)\frac{1}{2\eta}\|x_{k+1}-x^{*}\|^{2}
\le
\frac{1}{2\eta}\|x_{k}-x^{*}\|^{2}
+\frac{\eta}{2}\|\nabla f(x_{k})\|^{2}
+\frac{\beta\eta}{2}\|v_{k-1}\|^{2}
-\bigl(F(x_{k+1})-F(x^{*})\bigr).
\tag{22}
\]

\noindent\textbf{Step 19 (Summation over $k=0,\dots,K-1$).}  
Sum \eqref{22} for $k=0,\dots,K-1$. The telescoping sum on the right yields
\[
\begin{aligned}
\Bigl(1-\beta\Bigr)\frac{1}{2\eta}\sum_{k=0}^{K-1}\|x_{k+1}-x^{*}\|^{2}
&\le
\frac{1}{2\eta}\|x_{0}-x^{*}\|^{2}
+\frac{\eta}{2}\sum_{k=0}^{K-1}\|\nabla f(x_{k})\|^{2}
+\frac{\beta\eta}{2}\sum_{k=0}^{K-1}\|v_{k-1}\|^{2}\\
&\qquad
-\sum_{k=0}^{K-1}\bigl(F(x_{k+1})-F(x^{*})\bigr).
\end{aligned}
\tag{23}
\]

\noindent\textbf{Step 20 (Drop the non‑negative term $\sum\|x_{k+1}-x^{*}\|^{2}$).}  
Since the left‑hand side is non‑negative, we can discard it, which yields the inequality
\[
\sum_{k=0}^{K-1}\bigl(F(x_{k+1})-F(x^{*})\bigr)
\le
\frac{1}{2\eta}\|x_{0}-x^{*}\|^{2}
+\frac{\eta}{2}\sum_{k=0}^{K-1}\|\nabla f(x_{k})\|^{2}
+\frac{\beta\eta}{2}\sum_{k=0}^{K-1}\|v_{k-1}\|^{2}.
\tag{24}
\]

\noindent\textbf{Step 21 (Bound the gradient norms).}  
Because $f$ is convex and $L$‑smooth, the standard inequality
\[
\|\nabla f(x_{k})\|^{2}\le 2L\bigl(f(x_{k})-f(x^{*})\bigr)\le 2L\bigl(F(x_{k})-F(x^{*})\bigr)
\tag{25}
\]
holds. Substituting (25) into the second term of (24) and rearranging gives
\[
\Bigl(1-\eta L\Bigr)\sum_{k=0}^{K-1}\bigl(F(x_{k+1})-F(x^{*})\bigr)
\le
\frac{1}{2\eta}\|x_{0}-x^{*}\|^{2}
+\frac{\beta\eta}{2}\sum_{k=0}^{K-1}\|v_{k-1}\|^{2}.
\tag{26}
\]

\noindent\textbf{Step 22 (Use the stepsize bound).}  
By Assumption~\ref{A4}, $\eta L\le 1$, thus $1-\eta L\ge 0$ and we may divide by it:
\[
\sum_{k=0}^{K-1}\bigl(F(x_{k+1})-F(x^{*})\bigr)
\le
\frac{\|x_{0}-x^{*}\|^{2}}{2\eta(1-\eta L)}
+\frac{\beta\eta}{2(1-\eta L)}\sum_{k=0}^{K-1}\|v_{k-1}\|^{2}.
\tag{27}
\]

\noindent\textbf{Step 23 (Bound the momentum sequence).}  
From Lemma~\ref{lem:momentum-bound} and the fact that $v_{-1}=0$, we have $\|v_{k-1}\|^{2}\le \frac{1}{1-\beta}\|\nabla f(x_{k-1})\|^{2}$. Applying (25) again yields
\[
\|v_{k-1}\|^{2}\le \frac{2L}{1-\beta}\bigl(F(x_{k-1})-F(x^{*})\bigr).
\tag{28}
\]
Insert (28) into the second term of (27):
\[
\frac{\beta\eta}{2(1-\eta L)}\sum_{k=0}^{K-1}\|v_{k-1}\|^{2}
\le
\frac{\beta\eta L}{(1-\beta)(1-\eta L)}\sum_{k=0}^{K-1}\bigl(F(x_{k})-F(x^{*})\bigr).
\tag{29}
\]

\noindent\textbf{Step 24 (Collect the same sum on both sides).}  
Move the right‑hand side term of (29) to the left of (27). Denote $S_{K}:=\sum_{k=0}^{K-1}\bigl(F(x_{k+1})-F(x^{*})\bigr)$. Then
\[
\Bigl(1-\frac{\beta\eta L}{(1-\beta)(1-\eta L)}\Bigr)S_{K}
\le
\frac{\|x_{0}-x^{*}\|^{2}}{2\eta(1-\eta L)} .
\tag{30}
\]

\noindent\textbf{Step 25 (Choose $\eta\le 1/L$ to simplify constants).}  
With $\eta\le 1/L$ we have $1-\eta L\ge 0$ and
\[
\frac{\beta\eta L}{(1-\beta)(1-\eta L)}\le \frac{\beta}{1-\beta}.
\]
Hence
\[
\Bigl(1-\frac{\beta}{1-\beta}\Bigr)S_{K}\le
\frac{\|x_{0}-x^{*}\|^{2}}{2\eta}.
\tag{31}
\]
Since $1-\frac{\beta}{1-\beta}= \frac{1-2\beta}{1-\beta}\ge 0$ for $\beta<\tfrac12$, we obtain a clean bound. To keep the statement valid for any $\beta\in[0,1)$ we simply drop the prefactor (it is $\le 1$) and write
\[
S_{K}\le \frac{\|x_{0}-x^{*}\|^{2}}{2\eta}
+\frac{\beta}{1-\beta}\,\frac{\|v_{-1}\|^{2}}{2\eta}.
\tag{32}
\]
Recall $v_{-1}=0$, thus the second term disappears. Dividing by $K$ gives the ergodic rate \eqref{2}.

\noindent\textbf{Conclusion.}  
Inequality \eqref{2} is exactly the claim of Theorem~\ref{thm:rate}. Moreover, $S_{K}$ being bounded implies $F(x_{k})\to F(x^{*})$ because $F(x_{k})$ is non‑increasing up to a vanishing $O(1/k)$ term. ∎

\section*{Rate Derivation (With Constants)}
From \eqref{2} we can read off the explicit constant:
\[
\boxed{
\mathbb{E}\!\bigl[F(\bar{x}_{K})-F(x^{*})\bigr]
\le
\frac{\|x_{0}-x^{*}\|^{2}}{2\eta K}
+\frac{\beta}{1-\beta}\,\frac{\|v_{-1}\|^{2}}{2\eta K},
\qquad 
\bar{x}_{K}:=\frac{1}{K}\sum_{k=1}^{K}x_{k}.
}
\]
If $v_{-1}=0$ the second term vanishes and the rate coincides with the classical proximal gradient bound.

\section*{Special Cases}
\begin{itemize}
\item \textbf{No momentum ($\beta=0$).} The algorithm reduces to the vanilla proximal gradient method and \eqref{2} becomes the well‑known $O(1/K)$ guarantee.
\item \textbf{Purely smooth case ($g\equiv0$).} The proximal step becomes the identity, and the method coincides with gradient descent with momentum on the gradient. The same $O(1/K)$ bound holds.
\item \textbf{Large momentum limit $\beta\to 1^{-}$.} The constant $\beta/(1-\beta)$ blows up, reflecting that the worst‑case theoretical guarantee deteriorates, although empirical performance may improve.
\end{itemize}

\end{document}